\documentclass[8pt]{beamer}
\usepackage{tikz}
\usepackage{tikz-cd}
\usepackage{qtree}
\usepackage{amsfonts}
\usepackage{amssymb}
\usepackage{amsmath}
\usepackage{amscd}
\usepackage{amsthm}
% \usepackage{mathpartir}
\usepackage{bbm}
\usepackage[mathcal]{euscript}
% \usepackage[utf8]{inputenc}
\setbeamertemplate{navigation symbols}{}%remove navigation symbols
%\usepackage[T1]{fontenc}
\usepackage{lmodern} \title{Algebraic quantum field theory}
\author{Hans Halvorson}
\date[November 2013]{November 12, 2013} \usetheme{Copenhagen}
\begin{document}
\def\2{\mathcal}
\def\7{\mathbb}
\newenvironment{proenv}{\only{\setbeamercolor{local structure}{fg=green}}}{}
\newenvironment{conenv}{\only{\setbeamercolor{local structure}{fg=red}}}{}  
\newenvironment{neuenv}{\only{\setbeamercolor{local
      structure}{fg=yellow}}}{}


\newcommand{\3}{\mathfrak}

\begin{frame}[plain]
  \titlepage
\end{frame}

\begin{frame}{Haag-Kastler axioms}

  Let $\2M$ be Minkowski spacetime.  For each double cone $O$, we
  assume that $\2A (O)\subseteq \2A$ is a $C^*$-algebra.  The algebra
  $\2A$ is called the \textbf{quasilocal} algebra of observables.

  \bigskip \textbf{axiom 1 (isotony)} If $O_1\subseteq O_2$ then $\2A
  (O_1)\subseteq \2A (O_2)$.

  \bigskip \textbf{axiom 2 (microcausality)} If $O_1$ is spacelike
  separated from $O_2$ then $\2A (O_1)$ commutes with $\2A (O_2)$.

  \bigskip \textbf{axiom 3 (translation covariance)} There is a
  representation $x\mapsto \alpha _x$ of the translation group of
  $\2M$ in the group of automorphism of $\2A$ such that
$$ \alpha _x(\2A (O)) \: =\: \2A (O+x) .$$

\end{frame}

\begin{frame}{Haag-Kastler axioms}

  \textbf{positive energy vacuum representation:} Let $\omega$ be a
  translation invariant state of $\2A$, and let $(\2H ,\pi ,x\mapsto
  U(x),\Omega )$ be the covariant GNS representation induced by
  $\omega$.  Let $\2R (O)=\pi (\2A (O))''$.

  \bigskip \textbf{axiom 4 (additivity)} For each double cone $O$, the
  set $\{ \2R (O+x) :x\in \7R ^ 4 \}$ acts irreducibly on $\2H$.

  \bigskip \textbf{axiom 5 (spectrum condition)} The four-momentum
  $(P_0,P_1,P_2,P_3)$ has spectrum restricted to the forward light
  cone.

\end{frame}

\begin{frame}{first consequences of the axioms}


  \bigskip \textbf{Theorem (Reeh-Schlieder)} For any double cone $O$,
  the vacuum vector $\Omega$ is cyclic for $\2R (O)$.  That is, $\{
  A\Omega :A\in \2R (O) \}$ is dense in the Hilbert space $\2H$.

  \bigskip Proof uses an analyticity argument, invoking the spectrum
  condition.

\end{frame}
\begin{frame}{consequences of Reeh-Schlieder}

  \bigskip \hspace*{2em} $\hookrightarrow$ the vacuum state is
  entangled; is highly entropic in local regions

  \bigskip \hspace*{2em} $\hookrightarrow$ the vacuum vector is
  separating for localized operators: $A\Omega =0$ only if $A=0$

  \bigskip \hspace*{2em} $\hookrightarrow$ no localized number
  operators

\end{frame}

\begin{frame}{more consequences of the axioms}

  \textbf{Theorem (Kadison)} Local algebras are of infinite type.

  \bigskip \hspace*{2em} $\hookrightarrow$ corollary: for any two
  spacelike separated regions, a dense set of states violates Bell's
  inequality

  \bigskip \hspace*{2em} $\hookrightarrow$ all states are mixed (no
  normal states are pure)

\bigskip \hspace*{2em} $\hookrightarrow$ example: in the model of the
free Boson field, local algebras are type III$_1$.

\end{frame}


\begin{frame}{Bell's inequalities}

  Given $C^*$-algebras $\2A$ and $\2B$, such that $[A,B]= 0$ for all
  $A\in \2A$ and $B\in \2B$.

A \textcolor{blue}{Bell operator} is of the form
$$ R = \frac{1}{2}(A_1(B_1+B_2)+A_2(B_1-B_2)) ,$$
where $A_i\in \2A$ and $B_i\in \2B$ are self-adjoint with
$\| A_i\| \leq 1$ and $\| B_i \|\leq 1$.



\end{frame}

\begin{frame}{Bell's theorem}

  The correlations for measurements of $A_i$ and $B_i$ in state
  $\omega$ permit a \textcolor{blue}{local hidden variable model} just
  in case $|\omega (R)|\leq 1$.

\bigskip \bigskip \textcolor{blue}{Tsirelson bound:} All quantum states satisfy $|\omega (R)|\leq \sqrt{2}$.

\bigskip For $\2A = M_2\otimes I$ and $\2B=I\otimes M_2$, and
appropriate choice of spin operators $A_i$ and $B_i$, the singlet
state \textcolor{blue}{maximally violates} Bell's inequality.

\end{frame}

\begin{frame}{Nonlocality in QFT}

  \textbf{Theorem (Summers and Werner)} If $\2A$ and $\2B$ are the
  algebras associated with the Rindler wedges, then
  \textcolor{red}{every state} maximally violates Bell's inequality.

  \bigskip \bigskip \textbf{Theorem (Clifton and Halvorson)} If $\2A$
  and $\2B$ are algebras associated with any two spacelike separated
  regions, then the set of Bell-violating states is dense in the state
  space.



\end{frame}

\end{document}